\documentclass[11pt]{article}
\usepackage[margin=1in]{geometry}
\usepackage{amsmath,amssymb,amsthm}
\usepackage{graphicx}
\usepackage[hidelinks]{hyperref}
\setlength{\emergencystretch}{2em}

\newtheorem{theorem}{Theorem}
\newtheorem{lemma}{Lemma}
\newtheorem{corollary}{Corollary}

\newcommand{\C}{\mathbb C}
\newcommand{\supp}{\operatorname{supp}}

\title{A Trivial-Bundle Partial Answer to the C*-Algebra Ball-Covering Question}
\author{Run \texttt{fa\_banach\_001}, lane 19}
\date{June 17, 2026}

\begin{document}
\maketitle

\section*{Source Problem}

The source paper is Jinghao Huang, Karimbergen Kudaybergenov, and Rui Liu,
\emph{The ball-covering property of von Neumann algebras and noncommutative
symmetric spaces}, arXiv:2311.14257.  After proving that a commutative unital
\(C^*\)-algebra \(C(K)\) has BCP/UBCP exactly when its pure state space, i.e.
\(K\), has a countable \(\pi\)-basis, the authors ask for a noncommutative
version:

\begin{center}
\includegraphics[width=0.92\textwidth]{figures/open_problem_crop.png}
\end{center}

\section*{Statement}

Recall that a family of nonempty open sets in \(K\) is a \(\pi\)-basis if every
nonempty open set in \(K\) contains a member of the family.

\begin{theorem}
\label{thm:vector-c0}
Let \(K\) be a locally compact Hausdorff space and let \(X\) be a nonzero
Banach space.
\begin{enumerate}
    \item If \(C_0(K,X)\) has the ball-covering property, then \(K\) has a
    countable \(\pi\)-basis.
    \item If \(K\) has a countable \(\pi\)-basis and \(X\) has the
    ball-covering property, then \(C_0(K,X)\) has the ball-covering property.
    \item If \(K\) has a countable \(\pi\)-basis and \(X\) is separable, then
    \(C_0(K,X)\) has the uniform ball-covering property.
\end{enumerate}
\end{theorem}

\begin{corollary}
\label{cor:cstar}
Let \(A\) be a nonzero separable \(C^*\)-algebra.  Then \(C_0(K,A)\) has
BCP/UBCP if and only if \(K\) has a countable \(\pi\)-basis.
\end{corollary}

This gives a noncommutative extension of the commutative \(C_0(K)\) criterion
for the large class of trivial continuous \(C^*\)-bundles with separable fiber,
including \(C_0(K,M_n)\) for \(n\ge2\) and \(C_0(K,\mathcal K(H))\) for
separable \(H\).

After this packet was first drafted, a later literature check identified that
Liu--Liu--Lu--Zheng \cite{LLLZ} already state the broader Banach-valued
stability theorem for \(C_0(K,X)\), including strong and uniform versions.
Accordingly, the packet should be read as a C*-specialized self-contained
proof and application note for the source question, not as a novel theorem
beyond that vector-valued result.

\section*{Proof Intuition}

The obstruction is still topological.  If \(K\) has no countable \(\pi\)-basis,
then the near-norm sets of any countable family of candidate centers cannot
probe all open subsets of \(K\).  A vector-valued bump supported in an open set
missed by those near-norm sets is zero where each candidate center almost
attains its norm, so it is outside every admissible ball.

Conversely, a countable \(\pi\)-basis supplies countably many scalar bump
functions.  Near a point where a vector-valued function \(f\) nearly attains its
norm, choose a fiber ball center covering that fiber value and multiply it by a
small scalar bump supported where \(f\) changes little.  Outside the support the
distance is at most \(\|f\|\); inside the support, convexity controls the whole
transition from \(0\) to the fiber center.

\section*{Proof}

We use the standard equivalent form of BCP: the unit sphere of a normed space
\(Y\) is covered by countably many open balls \(B(y_n,\|y_n\|)\), \(y_n\ne0\).
Indeed, any open ball off the origin has radius at most the norm of its center,
and enlarging to radius \(\|y_n\|\) still keeps the origin out of the open ball.

\begin{lemma}[the base obstruction]
\label{lem:base-obstruction}
Let \(X\ne\{0\}\).  If \(K\) has no countable \(\pi\)-basis, then
\(C_0(K,X)\) fails BCP.
\end{lemma}

\begin{proof}
Let \((F_n)\) be any sequence in \(C_0(K,X)\).  We show that the balls
\(B(F_n,\|F_n\|)\) cannot cover the unit sphere.  For \(n,k\ge1\), set
\[
        U_{n,k}=\{t\in K:\|F_n(t)\|>\|F_n\|-1/k\}.
\]
These sets are open and nonempty whenever \(F_n\ne0\).  Since \(K\) has no
countable \(\pi\)-basis, there is a nonempty open set \(W\subset K\) such that
no \(U_{n,k}\) is contained in \(W\).  Choose \(h\in C_c(K)\) with
\(0\le h\le1\), \(\|h\|_\infty=1\), and \(\supp h\subset W\).  Choose
\(x_0\in X\) with \(\|x_0\|=1\), and put
\[
        G(t)=h(t)x_0 .
\]
Then \(G\in C_0(K,X)\) and \(\|G\|=1\).

Fix \(n\).  For every \(k\), choose \(t_{n,k}\in U_{n,k}\setminus W\).  Since
\(\supp h\subset W\), we have \(G(t_{n,k})=0\), hence
\[
        \|G-F_n\|
        \ge \|F_n(t_{n,k})\|
        > \|F_n\|-1/k .
\]
Letting \(k\to\infty\) gives \(\|G-F_n\|\ge\|F_n\|\).  Thus \(G\) is in none of
the open balls \(B(F_n,\|F_n\|)\).  Since the sequence \((F_n)\) was arbitrary,
\(C_0(K,X)\) fails BCP.
\end{proof}

\begin{lemma}[lifting a fiber BCP cover]
\label{lem:positive-bcp}
If \(K\) has a countable \(\pi\)-basis and \(X\) has BCP, then \(C_0(K,X)\)
has BCP.
\end{lemma}

\begin{proof}
Let \((V_j)\) be a countable \(\pi\)-basis of \(K\).  By local compactness and
regularity, for each \(j\) choose \(h_j\in C_c(K)\) such that
\[
        0\le h_j\le1,\qquad \|h_j\|_\infty=1,\qquad \supp h_j\subset V_j .
\]

Let \((a_m)\subset X\setminus\{0\}\) be BCP centers for \(X\), so that
\[
        S_X\subset \bigcup_m B(a_m,\|a_m\|).
\]
Discard centers whose balls miss \(S_X\), and replace the remaining centers by
\(y_m=4a_m\).  The cover property is preserved because if
\(\|x-a_m\|<\|a_m\|\), then
\[
        \|x-4a_m\|
        \le \|x-a_m\|+3\|a_m\|
        <4\|a_m\|=\|y_m\|.
\]
Also every retained \(a_m\) satisfies \(\|a_m\|>1/2\), so \(\|y_m\|>2\).

We claim that the countable family
\[
        G_{j,m}(t)=h_j(t)y_m
\]
is a BCP family for \(C_0(K,X)\).  Let \(f\in C_0(K,X)\) with \(\|f\|=1\).
Choose \(t_0\in K\) with \(\alpha:=\|f(t_0)\|>0\), and put
\(x=f(t_0)/\alpha\in S_X\).  Choose \(m\) such that
\(\|x-y_m\|<\|y_m\|\).  Since \(0\le\alpha\le1\),
\[
        \|\alpha x-y_m\|
        \le \alpha\|x-y_m\|+(1-\alpha)\|y_m\|
        <\|y_m\|.
\]
Moreover \(\|y_m\|>2\), so both endpoint distances
\(\|\alpha x\|\) and \(\|\alpha x-y_m\|\) are strictly smaller than
\(\|y_m\|\).  By convexity, the function
\[
        s\mapsto \|\alpha x-sy_m\|,\qquad 0\le s\le1,
\]
is bounded by the maximum of its endpoint values, and is therefore strictly
less than \(\|y_m\|\) on the whole interval.

Choose a neighbourhood \(U\) of \(t_0\) so small that
\[
        \|f(t)-f(t_0)\|
        <
        \|y_m\|-\sup_{0\le s\le1}\|\alpha x-sy_m\|
        \qquad(t\in U).
\]
Pick \(j\) with \(V_j\subset U\).  If \(t\notin\supp h_j\), then
\[
        \|f(t)-G_{j,m}(t)\|=\|f(t)\|\le1<\|y_m\|.
\]
If \(t\in\supp h_j\), then \(t\in U\) and
\[
        \|f(t)-G_{j,m}(t)\|
        \le
        \|f(t)-f(t_0)\|+\|\alpha x-h_j(t)y_m\|
        < \|y_m\|.
\]
Thus \(\|f-G_{j,m}\|<\|G_{j,m}\|\).  This covers the unit sphere.
\end{proof}

\begin{lemma}[separable fibers give UBCP]
\label{lem:separable-ubcp}
If \(K\) has a countable \(\pi\)-basis and \(X\) is separable, then
\(C_0(K,X)\) has UBCP.
\end{lemma}

\begin{proof}
Let \((u_m)\) be a countable \(1/16\)-net in \(S_X\), and let \((h_j)\) be the
scalar bumps from the preceding proof.  Consider the centers
\[
        G_{j,m}(t)=4h_j(t)u_m
\]
and use the fixed radius \(7/2\).  These balls are uniformly bounded and each
is disjoint from the open ball \(B(0,1/2)\), since \(\|G_{j,m}\|=4\).

Let \(f\in S_{C_0(K,X)}\).  Choose \(t_0\) with
\(\alpha=\|f(t_0)\|>3/4\), set \(x=f(t_0)/\alpha\), and choose \(m\) with
\(\|x-u_m\|<1/16\).  Then
\[
        \|\alpha x-4u_m\|
        \le \alpha\|x-u_m\|+(4-\alpha)
        < \frac1{16}+\frac{13}{4}<\frac72 .
\]
Also \(\|\alpha x\|<7/2\).  By convexity,
\[
        \sup_{0\le s\le1}\|\alpha x-4s u_m\|<7/2.
\]
Choose a neighbourhood \(U\) of \(t_0\) on which \(f\) varies by less than the
positive margin between \(7/2\) and the displayed supremum, and then choose
\(j\) with \(V_j\subset U\).  The same outside-support and inside-support
estimates as in Lemma~\ref{lem:positive-bcp} give
\[
        \|f-G_{j,m}\|<7/2 .
\]
Thus the unit sphere is covered by countably many balls of fixed radius
\(7/2\), all avoiding \(B(0,1/2)\).  This is UBCP.
\end{proof}

Theorem~\ref{thm:vector-c0} follows from Lemmas~\ref{lem:base-obstruction},
\ref{lem:positive-bcp}, and \ref{lem:separable-ubcp}.  Corollary~\ref{cor:cstar}
follows because every separable Banach space is a separable fiber in
Lemma~\ref{lem:separable-ubcp}, and Lemma~\ref{lem:base-obstruction} gives the
reverse implication for every nonzero fiber.

\section*{Scope and Limitations}

This is not a full characterization of arbitrary \(C^*\)-algebras with
BCP/UBCP.  It covers the important and natural class of trivial continuous
fields \(C_0(K,A)\) and shows that, in this class, the commutative base-space
condition is exactly the right obstruction.  The argument does not address
general nontrivial continuous fields, extensions, quotients, or arbitrary
nonseparable fibers.

\section*{Novelty and Search Bounds}

The run indexes were searched for arXiv:2311.14257, ball-covering property,
\(C^*\)-algebra, pure state space, \(\pi\)-basis, BCP, and UBCP.  No existing
packet for this target was found.  Web searches on June 17, 2026 found no
later direct characterization of all \(C^*\)-algebras with BCP/UBCP.  They did
identify arXiv:2111.04921 \cite{LLLZ}, which already proves the broader
Banach-valued \(C_0(K,X)\) stability theorem underlying this trivial-bundle
subcase, and arXiv:2507.02261 on approximation criteria for operator spaces.

\section*{Human Review Focus}

The main checks are:
\begin{itemize}
    \item in Lemma~\ref{lem:base-obstruction}, the use of near-norm open sets
    to turn failure of countable \(\pi\)-basis into a vector-valued bump
    escaping every countable ball cover;
    \item in Lemma~\ref{lem:positive-bcp}, the convexity estimate through the
    transition region of the scalar bump;
    \item in Lemma~\ref{lem:separable-ubcp}, the fixed constants \(4\), \(7/2\),
    and \(1/2\) used to obtain UBCP.
\end{itemize}

\section*{References}

\begin{thebibliography}{9}

\bibitem{HKL}
Jinghao Huang, Karimbergen Kudaybergenov, and Rui Liu,
\emph{The ball-covering property of von Neumann algebras and noncommutative
symmetric spaces},
arXiv:2311.14257, 2023.

\bibitem{LLLZ}
Minzeng Liu, Rui Liu, Jimeng Lu, and Bentuo Zheng,
\emph{Ball covering property from commutative function spaces to
non-commutative spaces of operators},
arXiv:2111.04921, 2021.

\end{thebibliography}

\end{document}
